Bu calismada, sentetik veri ureten bir toplu tasima telemetri sistemi gelistirilmis ve bu verinin MQTT tabanli bir akis ile AWS uzerinde gercek zamanli olarak islenmesi saglanmistir. Projede otobusler, belirli araliklarla konum, hiz, hat bilgisi ve kart basimina dayali binis sayisi ureten sanal cihazlar olarak modellenmistir. Uretilen telemetri verileri AWS IoT Core uzerinden sisteme alinmis, IoT kurali araciligiyla Amazon Kinesis Data Streams katmanina aktarilmis ve AWS Lambda fonksiyonu ile zenginlestirilerek Amazon DynamoDB tablolarina yazilmistir. Islenen veriler FastAPI tabanli bir okuma katmani ile sunulmus ve Streamlit tabanli bir dashboard uzerinden canli olarak gorsellestirilmistir.

Calismada ham veri ile turetilmis veri acik bicimde ayrilmistir. Konum, hiz ve binis sayisi simulator tarafinda uretilirken; tahmini varis suresi, tahmini inis miktari, doluluk seviyesi ve gecikme bilgisi bulut tarafinda hesaplanmistir. Bu yaklasim, projeyi yalnizca gorsel bir simulasyon olmaktan cikarmis ve gercek zamanli veri isleme mimarisinin uctan uca hayata gecirildigi bir uygulamaya donusturmustur. Elde edilen sonuc, AWS servisleri kullanilarak kurulan MQTT tabanli veri hatti ile sentetik otobus telemetrisinin basarili bicimde islenebildigini, depolanabildigini ve canli izleme amacli bir arayuze aktarilabildigini gostermistir.
